\section{Results}

\subsection{Segmentation Quality}

Precision, recall, and F1-score refer to the report samples described in the experimental setup. All quality metrics below are expressed as percentages.

\subsubsection{BRIK}

\begin{table}[H]
  \centering
  \caption{BRIK segmentation quality with the voxel-sampler-based pipeline.}
  \label{tab:brik_voxel}
  \small
  \begin{tabular}{lrrrrr}
    \toprule
    Class & Precision & Recall & F1 & IoU & Support \\
    \midrule
    Railway tracks & 93 & 34 & 50 & $33.46$ & 66\,009 \\
    Ground & 52 & 95 & 67 & $50.91$ & 54\,003 \\
    Masts & 69 & 38 & 49 & $32.57$ & 24\,683 \\
    Support devices & 57 & 75 & 65 & $47.99$ & 25\,433 \\
    Overhead lines & 84 & 84 & 84 & $72.40$ & 18\,647 \\
    Fences & 99 & 26 & 41 & $25.85$ & 23\,513 \\
    Poles & 83 & 55 & 66 & $49.53$ & 25\,297 \\
    Vegetation & 59 & 96 & 73 & $57.37$ & 66\,001 \\
    Buildings & 75 & 79 & 77 & $62.93$ & 26\,824 \\
    Other objects & 82 & 39 & 53 & $35.73$ & 33\,139 \\
    \bottomrule
  \end{tabular}
\end{table}

The BRIK report for the voxel-sampler-based pipeline gives an accuracy of 64.43\% and mIoU of 46.87\%. Railway tracks, ground, and vegetation account for nearly half of the report sample. Fence predictions are predominantly correct but miss much of the reference class, whereas ground and vegetation include many false positives.

\begin{table}[H]
  \centering
  \caption{BRIK segmentation quality with the probability-sampler-based pipeline.}
  \label{tab:brik_probability}
  \small
  \begin{tabular}{lrrrrr}
    \toprule
    Class & Precision & Recall & F1 & IoU & Support \\
    \midrule
    Railway tracks & 99 & 56 & 72 & $5.64$ & 66\,009 \\
    Ground & 44 & 98 & 61 & $82.29$ & 54\,003 \\
    Masts & 32 & 59 & 41 & $1.44$ & 24\,683 \\
    Support devices & 71 & 13 & 23 & $5.86$ & 25\,433 \\
    Overhead lines & 85 & 88 & 87 & $68.28$ & 18\,647 \\
    Fences & 87 & 24 & 38 & $12.60$ & 23\,513 \\
    Poles & 75 & 1 & 1 & $<0.001$ & 25\,297 \\
    Vegetation & 53 & 73 & 61 & $76.34$ & 66\,001 \\
    Buildings & 41 & 55 & 47 & $26.32$ & 26\,824 \\
    Other objects & 43 & 26 & 32 & $7.09$ & 33\,139 \\
    \bottomrule
  \end{tabular}
\end{table}

For the probability-sampler-based pipeline on BRIK, accuracy is 82.30\%, balanced accuracy is 45.06\%, and mIoU is 28.59\%. Both BRIK tables use identical per-class support. Pole recall is close to zero in a~sample containing tens of thousands of reference pole points.

\subsubsection{DALES}

\begin{table}[H]
  \centering
  \caption{DALES segmentation quality with the probability-sampler-based pipeline.}
  \label{tab:dales_probability}
  \small
  \begin{tabular}{lrrrrr}
    \toprule
    Class & Precision & Recall & F1 & IoU & Support \\
    \midrule
    Ground & 86 & 99 & 92 & $92.82$ & 76\,338 \\
    Vegetation & 70 & 94 & 80 & $91.87$ & 76\,337 \\
    Cars & 79 & 91 & 85 & $78.73$ & 76\,336 \\
    Trucks & 93 & 37 & 53 & $19.85$ & 56\,159 \\
    Power lines & 94 & 93 & 94 & $92.42$ & 52\,528 \\
    Fences & 97 & 77 & 86 & $52.31$ & 76\,335 \\
    Poles & 95 & 74 & 83 & $55.02$ & 59\,635 \\
    Buildings & 74 & 89 & 81 & $82.54$ & 76\,332 \\
    \bottomrule
  \end{tabular}
\end{table}

The DALES report for the probability-sampler-based pipeline gives accuracy of 94.85\%, balanced accuracy of 81.29\%, and mIoU of 70.70\%. The report sample contains fewer truck, pole, and power-line points than points from the other classes. Power-line predictions combine high precision and recall, while truck errors are dominated by missed reference points rather than false detections.

\begin{table}[H]
  \centering
  \caption{DALES segmentation quality with the voxel-sampler-based pipeline.}
  \label{tab:dales_voxel}
  \small
  \begin{tabular}{lrrrrr}
    \toprule
    Class & Precision & Recall & F1 & IoU & Support \\
    \midrule
    Ground & 77 & 96 & 86 & $74.88$ & 76\,338 \\
    Vegetation & 59 & 96 & 73 & $57.84$ & 76\,337 \\
    Cars & 77 & 90 & 83 & $70.52$ & 76\,336 \\
    Trucks & 94 & 37 & 53 & $36.50$ & 56\,159 \\
    Power lines & 96 & 74 & 84 & $71.77$ & 52\,528 \\
    Fences & 99 & 65 & 79 & $64.62$ & 76\,335 \\
    Poles & 96 & 75 & 84 & $72.79$ & 59\,635 \\
    Buildings & 87 & 95 & 91 & $82.94$ & 76\,332 \\
    \bottomrule
  \end{tabular}
\end{table}

The DALES report for the voxel-sampler-based pipeline gives an accuracy of 80.27\% and mIoU of 66.48\%. Truck predictions show the largest precision--recall imbalance: false detections are uncommon, but many reference truck points are missed.

Both DALES report tables contain 550\,000 points with identical per-class support. The largest F1-score differences concern power lines and buildings: the probability-sampler-based pipeline performs better on the former, while the voxel-sampler-based pipeline performs better on the latter. Truck recall remains similarly limited in both reports.

\subsection{Performance Evaluation}

The benchmarks compare the same model used with the probability-sampler-based and voxel-sampler-based pipelines on 10 DALES scenes. Figures~\ref{fig:processing_time} and~\ref{fig:processing_stages} report mean values across these 10 scenes.

\begin{figure}[H]
  \centering
  \comparisonplot{time}{Total processing,Inference,CPU time}{Time [s]}{800}{1}
  \caption{Mean elapsed processing time and accumulated CPU time across 10 DALES scenes.}
  \label{fig:processing_time}
\end{figure}

Mean total elapsed time was 527.70~s for the probability-sampler-based pipeline and 238.44~s for the voxel-sampler-based pipeline, a~54.8\% reduction. Mean end-to-end throughput was approximately 22.3 thousand and 49.6 thousand input points per second, respectively. Mean accumulated CPU time was 678.92~s and 638.96~s, respectively.

\begin{figure}[H]
  \centering
  \comparisonplot{inference}{Model prediction,Reconstruction,Other inference}{Time [s]}{400}{1}
  \caption{Mean inference times by stage across 10 DALES scenes. ``Other inference'' is total inference time minus model prediction and recorded reconstruction time. Reconstruction was not recorded for the probability-sampler-based pipeline (n/r).}
  \label{fig:processing_stages}
\end{figure}

Mean model prediction time was 339.86~s for the probability-sampler-based pipeline and 13.50~s for the voxel-sampler-based pipeline. For the voxel-sampler-based pipeline, most inference time falls outside model prediction and reconstruction; mean recorded reconstruction time was 26.57~s. The logs for the probability-sampler-based pipeline contain no separate reconstruction stage; any reconstruction work is included in ``Other inference''.

\begin{figure}[H]
  \centering
  \comparisonplot{memory}{Baseline RSS,Peak PSS,Peak RSS,RSS increase}{Memory [GiB]}{75}{1}
  \caption{Mean process-tree RAM measurements across 10 DALES scenes.}
  \label{fig:memory_consumption}
\end{figure}

Mean peak PSS was 2.72~GiB for the probability-sampler-based pipeline and 5.04~GiB for the voxel-sampler-based pipeline. Mean peak RSS was 2.81~GiB and 64.42~GiB, respectively. Process-tree RSS can count shared pages in multiple worker processes; it should not be interpreted as distinct physical memory use.

\begin{figure}[H]
  \centering
  \resourcetimeline{probability}{probabilitycolor}{\ProbabilityTimeMax}{\ProbabilityCpuMax}{\ProbabilityRamMax}
  \caption{Mean CPU and RAM usage for the probability-sampler-based pipeline across 10 DALES scenes. CPU usage of 100\% corresponds to one logical CPU.}
  \label{fig:resource_timeline}
\end{figure}

\begin{figure}[H]
  \centering
  \resourcetimeline{voxel}{voxelcolor}{\VoxelTimeMax}{\VoxelCpuMax}{\VoxelRamMax}
  \caption{Mean CPU and RAM usage for the voxel-sampler-based pipeline across 10 DALES scenes. CPU usage of 100\% corresponds to one logical CPU.}
  \label{fig:resource_timeline_voxel}
\end{figure}

Figures~\ref{fig:resource_timeline} and~\ref{fig:resource_timeline_voxel} show mean resource traces, with CPU usage on the left axis and memory on the right. Traces are aligned by the fraction of each run completed before averaging; the time axis is scaled to the mean duration.

\begin{figure}[H]
  \centering
  \voxelsizeplot
  \caption{DALES accuracy and prediction coverage for voxel sizes from 0.1 to 0.5~m on three test clouds. Accuracy is evaluated on labeled points with returned predictions; the evaluated point sets differ between pipelines and voxel sizes. Solid lines show accuracy and dashed lines show coverage. Coverage is the number of returned predictions divided by the sparse input point count.}
  \label{fig:dales_voxel_size}
\end{figure}

\subsection{Qualitative Analysis}

BRIK provides a~supplementary comparison on strongly class-imbalanced railway scans with generally low, spatially variable point density and large gaps in coverage. Even after class balancing, the three largest classes---railway tracks, ground, and vegetation---together have about ten times the report support of overhead lines (Tables~\ref{tab:brik_voxel} and~\ref{tab:brik_probability}). The largest F1-score gains favor the voxel-sampler-based pipeline for poles and support devices, whereas railway tracks favor the probability-sampler-based pipeline. Overhead-line F1 is similar for both pipelines and exceeds 80\%. Fence predictions have high precision, but both pipelines miss roughly three-quarters of the reference fence points. The voxel-sampler-based pipeline has higher mIoU despite lower overall accuracy.

DALES provides the main comparison between the voxel-sampler-based and probability-sampler-based pipelines (Tables~\ref{tab:dales_probability} and~\ref{tab:dales_voxel}). Class support is more even in these report samples: the largest class has roughly one and a~half times as many points as the smallest. The probability-sampler-based pipeline leads in mIoU by about four percentage points. Power-line and building F1-scores differ by ten percentage points in opposite directions: power lines favor the probability-sampler-based pipeline, while buildings favor the voxel-sampler-based pipeline. The latter also has higher building precision and recall. Trucks remain the weakest class in both reports: nearly two-thirds of their reference points are missed, despite few false-positive truck predictions.

Across the displayed voxel-size range of 0.1--0.5~m, the accuracy loss for the probability-sampler-based pipeline is more than twice that for the voxel-sampler-based pipeline (Figure~\ref{fig:dales_voxel_size}). The voxel-sampler-based pipeline's accuracy lead grows from roughly two to eight percentage points. Its prediction coverage remains complete, while that of the probability-sampler-based pipeline increases to nearly complete coverage at 0.5~m.
